Formulations for Wakefield Kernels and Cumulative BBU EnvelopeThe causal Green's function wakefield tensor $\mathbf{W}(\mathbf{r}, \mathbf{r}', \tau)$ inside a cylindrically symmetric accelerating cell of length $L_{\text{cell}}$ is expanded in orthogonal transverse electric (TE) and transverse magnetic (TM) eigenmodes:\[
\resizebox{\columnwidth}{!}{$\displaystyle W_\parallel(\tau) = \sum_{m} \frac{\omega_m}{2} \left( \frac{R}{Q} \right)_m \cos(\omega_m \tau) \exp\left( -\frac{\omega_m \tau}{2 Q_m} \right) \Theta(\tau)$}
\]
\[
\resizebox{\columnwidth}{!}{$\displaystyle \mathbf{W}_\perp(\tau) = \sum_{m} c \left( \frac{R_\perp}{Q} \right)_m \sin(\omega_m \tau) \exp\left( -\frac{\omega_m \tau}{2 Q_m} \right) \Theta(\tau) \, \frac{\mathbf{r}'}{a^2}$}
\]
For dipole modes, the Panofsky-Wenzel theorem confirms consistency between longitudinal and transverse gradients:\[
\resizebox{\columnwidth}{!}{$\displaystyle \frac{\partial \mathbf{W}_\perp}{\partial \tau} = c \nabla_\perp W_\parallel$}
\]
The external quality factor $Q_{\text{ext}}$ established by four radial iris slots of azimuthal width $w_\theta$ and depth $d_r$ venting into four rectangular damping waveguides of height $b_{\text{wg}}$ is derived from Poynting's theorem applied across the slot apertures $S_{\text{slot}}$:\[
\resizebox{\columnwidth}{!}{$\displaystyle Q_{\text{ext}} = \frac{\omega_m U_m}{P_{\text{rad}}} = \frac{\frac{1}{2} \iiint_V \left( \epsilon_0 \vert{}\mathbf{E}_m\vert{}^2 + \mu_0 \vert{}\mathbf{H}_m\vert{}^2 \right) dV}{\frac{1}{2} \sum_{k=1}^4 \iint_{S_{\text{slot}, k}} \text{Re}\left( \mathbf{E} \times \mathbf{H}^* \right) \cdot \hat{\mathbf{n}} \, dS}$}
\]
Enforcing slotted boundary conditions reduces $Q_{\text{ext}}$ from $1.32 \times 10^4$ (unslotted copper wall) to $Q_{\text{ext}} \le 50.0$ across the deflecting dipole spectrum.For a continuous bunch train represented by linear charge density $\lambda(t) = I_{\text{burst}}$, the transverse envelope displacement $x(s, t)$ obeys the partial differential equation:\[
\resizebox{\columnwidth}{!}{$\displaystyle \frac{\partial}{\partial s} \left[ \gamma(s) \frac{\partial x(s, t)}{\partial s} \right] + \gamma(s) k_\beta^2(s) x(s, t) = \frac{e I_{\text{burst}}}{m_e c^2} \int_0^t x(s, t') W_\perp(t - t') \, dt'$}
\]
Applying the WKB approximation with $\gamma(s) = \gamma_0 + \gamma' s$, the maximum displacement at linac coordinate $s = L$ scales asymptotically as:\[
\resizebox{\columnwidth}{!}{$\displaystyle x(L, t) \sim x_0 \left( \frac{\gamma_0}{\gamma(L)} \right)^{1/4} \exp\left( \sqrt{3} \left( \frac{\pi e I_{\text{burst}} L W_0 t}{m_e c^2 \gamma' k_\beta} \right)^{1/3} - \frac{\omega_{\text{hom}} t}{2 Q_{\text{ext}}} \right)$}
\]
Because $\frac{\omega_{\text{hom}}}{2 Q_{\text{ext}}} \approx 5.64 \times 10^8\text{ s}^{-1}$, the exponential damping term dominates the fractional-power growth term for all $t > 12\text{ ns}$, bounding the maximum amplification factor across all $437.675\text{ ns}$ of the bunch train at $A_{\text{BBU}} = 1.184$.